% Options for packages loaded elsewhere
\PassOptionsToPackage{unicode}{hyperref}
\PassOptionsToPackage{hyphens}{url}
%
\documentclass[
]{article}
\usepackage{amsmath,amssymb}
\usepackage{lmodern}
\usepackage{iftex}
\ifPDFTeX
  \usepackage[T1]{fontenc}
  \usepackage[utf8]{inputenc}
  \usepackage{textcomp} % provide euro and other symbols
\else % if luatex or xetex
  \usepackage{unicode-math}
  \defaultfontfeatures{Scale=MatchLowercase}
  \defaultfontfeatures[\rmfamily]{Ligatures=TeX,Scale=1}
\fi
% Use upquote if available, for straight quotes in verbatim environments
\IfFileExists{upquote.sty}{\usepackage{upquote}}{}
\IfFileExists{microtype.sty}{% use microtype if available
  \usepackage[]{microtype}
  \UseMicrotypeSet[protrusion]{basicmath} % disable protrusion for tt fonts
}{}
\makeatletter
\@ifundefined{KOMAClassName}{% if non-KOMA class
  \IfFileExists{parskip.sty}{%
    \usepackage{parskip}
  }{% else
    \setlength{\parindent}{0pt}
    \setlength{\parskip}{6pt plus 2pt minus 1pt}}
}{% if KOMA class
  \KOMAoptions{parskip=half}}
\makeatother
\usepackage{xcolor}
\usepackage[margin=1in]{geometry}
\usepackage{graphicx}
\makeatletter
\def\maxwidth{\ifdim\Gin@nat@width>\linewidth\linewidth\else\Gin@nat@width\fi}
\def\maxheight{\ifdim\Gin@nat@height>\textheight\textheight\else\Gin@nat@height\fi}
\makeatother
% Scale images if necessary, so that they will not overflow the page
% margins by default, and it is still possible to overwrite the defaults
% using explicit options in \includegraphics[width, height, ...]{}
\setkeys{Gin}{width=\maxwidth,height=\maxheight,keepaspectratio}
% Set default figure placement to htbp
\makeatletter
\def\fps@figure{htbp}
\makeatother
\setlength{\emergencystretch}{3em} % prevent overfull lines
\providecommand{\tightlist}{%
  \setlength{\itemsep}{0pt}\setlength{\parskip}{0pt}}
\setcounter{secnumdepth}{-\maxdimen} % remove section numbering
\newlength{\cslhangindent}
\setlength{\cslhangindent}{1.5em}
\newlength{\csllabelwidth}
\setlength{\csllabelwidth}{3em}
\newlength{\cslentryspacingunit} % times entry-spacing
\setlength{\cslentryspacingunit}{\parskip}
\newenvironment{CSLReferences}[2] % #1 hanging-ident, #2 entry spacing
 {% don't indent paragraphs
  \setlength{\parindent}{0pt}
  % turn on hanging indent if param 1 is 1
  \ifodd #1
  \let\oldpar\par
  \def\par{\hangindent=\cslhangindent\oldpar}
  \fi
  % set entry spacing
  \setlength{\parskip}{#2\cslentryspacingunit}
 }%
 {}
\usepackage{calc}
\newcommand{\CSLBlock}[1]{#1\hfill\break}
\newcommand{\CSLLeftMargin}[1]{\parbox[t]{\csllabelwidth}{#1}}
\newcommand{\CSLRightInline}[1]{\parbox[t]{\linewidth - \csllabelwidth}{#1}\break}
\newcommand{\CSLIndent}[1]{\hspace{\cslhangindent}#1}
\ifLuaTeX
  \usepackage{selnolig}  % disable illegal ligatures
\fi
\IfFileExists{bookmark.sty}{\usepackage{bookmark}}{\usepackage{hyperref}}
\IfFileExists{xurl.sty}{\usepackage{xurl}}{} % add URL line breaks if available
\urlstyle{same} % disable monospaced font for URLs
\hypersetup{
  pdftitle={Portrait des collaborations entre les etudiants du cours BIO500},
  pdfkeywords={Collaborations, Reseau ecologique, Travaux d'equipe},
  hidelinks,
  pdfcreator={LaTeX via pandoc}}

\title{Portrait des collaborations entre les etudiants du cours BIO500}
\author{true \and true \and true \and true}
\date{2023-04-20}

\begin{document}
\maketitle

\textbf{Resume}

La collaboration est une forme de travail où l'énergie est mise à profit
à l'aide de plusieurs personnes permettant l'atteinte d'un but commun
qui sera favorable pour tous. C'est cette forme de travail qui sera
décortiquée dans ce travail à l'aide des différentes collaborations que
les élèves en Baccalauréat en écologie de l'Université de Sherbrooke
dans le cadre du cours BIO500 ont eu à avoir tout au long de leur
cheminement. Ce travail vise à comprendre à quoi ressemble le réseau de
collaboration. Nous nous sommes intéressés à savoir si la distribution
du nombre total de collaborations de chaque étudiant suivait une
tendance particulière. De plus, nous avons regardé si certains cours
possédaient un plus grand nombre (ou le contraire) de collaborations
recensées. Il a été démontré à l'aide des données obtenues que plusieurs
étudiants réalisent plus de collaborations que d'autres, mais la
majorité s'en tiennent en bas de 50 collaborations au total.

\textbf{Introduction}

La collaboration est une forme d'aide où plusieurs personnes mettent en
oeuvre ces compétences pour l'atteinte d'un but commun précis, et ce
dans plusieurs sphères. En écologie, il a été démontré à plusieurs
reprises que la collaboration à son lot d'avantages concrets dans
plusieurs sphères soit dans le milieu animal, végétal et même chez
l'humain. La recherche en écologie, nécessite l'apport de plusieurs
connaissances adjacentes afin d'obtenir des résultats plus fiables vu la
complexité de la discipline (Goring et al. 2014). C'est principalement
cette interaction qui sera décortiquée dans le cadre du cours Méthodes
en écologie computationnelle enseigné aux étudiants des différentes
branches du baccalauréat en biologie de l'Université de Sherbrooke. En
effet, une collecte de données a été réalisée afin d'évaluer la
collaboration de chacun des étudiants présents dans ce cours avec
d'autres élèves tout au long de leur cheminement académique. Les trois
questions auxquelles notre équipe a décidé de s'attarder sont : est-ce
que le nombre total de collaborations pour chaque étudiant suit une
tendance particulière? Quels cours ont un plus nombre de collaborations?
Et finalement, à quoi ressemble le réseau de collaboration total ?

\textbf{Méthode et resultats}

\textbf{Collecte de la base de données}

En généralité, une collecte de données a été réalisée en sondant
l'entièreté des élèves présents dans le cours méthodes computationnelles
(BIO500) au Baccalauréat en écologie en hiver 2023. Chaque équipe devait
dument remplir 3 fichiers Excel avec différentes informations sur les
étudiants, les cours et les collaborations. Pour chaque étudiant
(premier fichier), on s'intéressait à connaitre sa région
administrative, sa présence ou non dans le régime coopératif, sa
formation préalable, l'année de début de son baccalauréat et son
programme d'étude. Pour chaque cours impliqué dans une collaboration
entre étudiants (deuxième fichier), on voulait connaitre le sigle de ce
cours, s'il était optionnel ou non et son nombre de crédits associé.
Pour chaque collaboration entre deux étudiants (troisième fichier), on
s'intéressait à connaitre le prénom et le nom de ces deux étudiants, le
sigle du cours dans lequel ils ont collaboré et la session pendant
laquelle cette collaboration a eu lieu. Ces fichiers Excel ont
finalement été transformés en fichiers `.csv' .

À l'aide du logiciel informatique R, nous avons compilé et fusionné tous
les fichiers .csv de données récoltées de chaque équipe dans trois
fichiers distincts, soit ``étudiants'', ``collaboration'' et ``cours''.
Pour chacun de ces fichiers R, un nettoyage complet et en profondeur des
données a été complété afin d'éliminer toutes les incohérences, les
erreurs de typographie et les doublons qui ont été insérés par
inattention lors de la saisie des informations par les différentes
équipes. Ce nettoyage avait pour objectif de rendre les bases de données
uniformes pour pouvoir s'en servir plus tard lors d'analyses.

\includegraphics{images/Reseau_de_collaboration.png}

Figure 1: Réseau des collaborations entre les différents étudiants du
cours BIO500 au Baccaulauréat en écologie.

Le but de cette figure 1 était de représenter le réseau de collaboration
des étudiants. Les différents noeuds représentent l'ensemble des
étudiants en plusieurs couleurs. Plus le noeud illustré est grand, plus
cet étudiant a réalisé un grand nombre de collaborations durant tout le
cheminement de ses études. Les lignes représentent chaque collaboration
entre ces étudiants. La couleur des points a été attribuée selon le
nombre de collaborations, soit bleu foncé pour plus de collaborations et
orange brulée pour un étudiant qui a eu le moins de collaborations. Ces
attributions de couleurs sont générées à l'aide de rangs préalablement
classée par le nombre de collaborations total additionné.

\hfill\break
\includegraphics{images/Histogramme_de_collaboration_par_cours.png}Figure
2 : Histogramme du nombre de collaborations selon les différents cours
présents dans le cheminement du Baccalauréat en écologie.

Le but de la deuxième figure était de représenter le nombre total de
collaborations pour chaque étudiant. Nous avons fait une requête SQL qui
nous a permis de créer un data.frame avec le nom et prénom de chacun des
étudiants impliqués dans le réseau. La première colonne recueille
l'ensemble des cours au Baccalauréat et le nombre de collaborations qui
ont été faites. À partir de cette requête, un histogramme a été
construit à partir de la fonction hist() et cette figure a été
sauvegardée dans un fichier `.png'. La Figure 2 permet de répondre à la
question suivante: quels cours ont un plus nombre de collaborations?
Parmi les différents cours qui comportent des collaborations, TSB303 est
celui qui possède le plus grand nombre avec un recensement de 579
collaborations. Par la suite, BCM104 en compte 383, ZOO304 en compte
224, ECL527 en compte 214 et ainsi de suite.

\includegraphics{images/Histogramme_de_collaboration_entre_etudiants.png}\includegraphics{}Figure
3: Histogramme des liens de collaboration répartit selon le nombre
d'étudiants.

Le but de la troisième figure était de représenter le nombre total de
collaborations de l'ensemble des étudiants pour chaque cours. Nous avons
fait une requête SQL qui nous a permis de créer un data.frame avec le
sigle des cours dans la première colonne et le nombre de collaborations
qui a été effectué pour chacun dans la deuxième colonne. À partir de
cette requête, un histogramme a été construit avec la fonction hist() et
cette figure a été sauvegardée dans un fichier `.png'. La figure 3 nous
aide à répondre à la question suivante: est-ce que le nombre total de
collaborations pour chaque étudiant suit une tendance particulière? La
majorité des étudiants se retrouvent dans la première classe, soit entre
0 et 50 collaborations pour un nombre d'étudiants totalisant une somme
plus grand que 100. Donc, il est possible de comprendre que la majorité
des étudiants a réalisée un nombre inférieur à 50 collaborations sur
l'entièreté de leur cheminement scolaire.

\textbf{Discussion}

En se basant sur la Figure 1 du réseau de collaboration, on remarque que
les étudiants ont tendance à préserver les mêmes collaborateurs tout au
long de leur parcours universitaire. En effet, en prenant en compte que
ce sont un total de 156 étudiants qui font partie de la base de données,
seulement moins de 10 étudiants ont plus de 150 collaborations. D'après
(Cullen et al. 1999), il est souhaitable de faire diverses
collaborations, puisque ces collaborateurs peuvent en tirer des
avantages scientifiques puisqu'elles mettent à profit les forces de
chaque individu. Par la suite, selon la Figure 2, les deux cours ayant
le plus de collaborations sont TSB303 et BCM104, deux cours obligatoires
pour obtenir son diplôme dans une branche de la biologie, soit
l'écologie. Ces 2 cours se dissocient de la courbe obtenue avec les
autres cours du baccalauréat par le haut nombre de collaborations
comptabilisé. Ces interactions entre plusieurs individus auraient lieu
lorsque les bénéfices du partenariat dépassent les couts engendrés, ce
qui a probablement eu lieu pour ces 2 cours (Goring et al. 2014).
Finalement, en ce qui concerne la question portant sur la tendance du
nombre de collaborations par étudiant qui correspond à la Figure 3, on
peut constater qu'uniquement très peu d'étudiants sur un total de 156
étudiants font plus de 150 collaborations, ce qui représente une infirme
partie de la majorité. Selon (REYERS et al. 2010), le domaine de
l'écologie au niveau universitaire ainsi que les chercheurs en début de
carrière se concentrent sur une seule discipline, diminuant le niveau de
collaboration.

\hypertarget{bibliographie}{%
\section*{Bibliographie}\label{bibliographie}}
\addcontentsline{toc}{section}{Bibliographie}

\hfill\break

\hfill\break

\hypertarget{refs}{}
\begin{CSLReferences}{1}{0}
\leavevmode\vadjust pre{\hypertarget{ref-cullen1999}{}}%
Cullen, Peter W., Richard H. Norris, Vincent H. Resh, Trefor B.
Reynoldson, David M. Rosenberg, and Michael T. Barbour. 1999.
{``Collaboration in Scientific Research: A Critical Need for Freshwater
Ecology.''} \emph{Freshwater Biology} 42 (1): 131--42.
\url{https://doi.org/10.1046/j.1365-2427.1999.00447.x}.

\leavevmode\vadjust pre{\hypertarget{ref-goring2014improving}{}}%
Goring, Simon J, Kathleen C Weathers, Walter K Dodds, Patricia A
Soranno, Lynn C Sweet, Kendra S Cheruvelil, John S Kominoski, Janine
Rüegg, Alexandra M Thorn, and Ryan M Utz. 2014. {``Improving the Culture
of Interdisciplinary Collaboration in Ecology by Expanding Measures of
Success.''} \emph{Frontiers in Ecology and the Environment} 12 (1):
39--47.

\leavevmode\vadjust pre{\hypertarget{ref-reyers2010}{}}%
REYERS, BELINDA, DIRK J. ROUX, RICHARD M. COWLING, AIMEE E. GINSBURG,
JEANNE L. NEL, and PATRICK O'FARRELL. 2010. {``Conservation Planning as
a Transdisciplinary Process.''} \emph{Conservation Biology} 24 (4):
957--65. \url{https://doi.org/10.1111/j.1523-1739.2010.01497.x}.

\end{CSLReferences}

\end{document}
